\chapter{项目背景与目标}
\label{ch:problems}

本章说明项目的应用背景、立项依据与要解决的问题，并给出系统的目标定义。

\section{项目定位}
\label{sec:positioning}

本项目设计并实现一台面向工业巡检场景的自主移动机器人，目标是：代替人工在工厂、变电站、机房等场景中自主遍历全区域，自动读取并分析工业仪表读数，并将过程与结果实时呈现在管理终端上。为此，系统需同时具备三类能力：

\begin{enumerate}
\item 自主移动与全覆盖。在无人工遥控、无预先地图的条件下自主完成建图与定位（SLAM），规划一条不漏扫、低重复的巡检路径，并在行进中实时躲避障碍。
\item 视觉读数与分析。在移动中实时识别画面中的工业仪表，定位表盘关键点并换算出指针读数，判断是否越限告警，并能调用大模型对历史数据生成趋势分析报告。
\item 多机协同与统一展示。多台机器人在同一坐标系下分区协作、协同覆盖大区域；所有机器人的位姿、进度与识别结果实时汇聚到同一台平板终端。
\end{enumerate}

上述能力均由本项目自主实现，并运行在国产软硬件生态之上。

\section{工业巡检的需求与痛点}
\label{sec:inspection-need}

在电力、石化、冶金、轨道交通、数据中心等场景中，抄录现场大量指针式仪表（压力表、油温表、SF\textsubscript{6} 气体密度表等）是巡检作业中占用人力最多、也最易出错的环节。此类仪表无源可靠、抗干扰、防爆性好，短期内难以被数字传感器全面替代。传统人工巡检存在三方面问题：

\begin{itemize}
\item 效率低、成本高。据行业实践数据，设备密集区域人工完成全区域巡检约需 300~min，智能系统仅需约 120~min，效率相差 2.5 倍以上。
\item 安全风险大。巡检常涉及高压、高温、有毒、易燃易爆等危险环境，极端天气与夜间作业时风险更为突出。
\item 数据难以追溯。纸质记录形成信息孤岛，历史数据无法用于趋势研判与预测性维护。
\end{itemize}

因此，以机器人替代人工巡检已成为行业趋势。市场研究也印证了这一方向：多家机构预测全球巡检机器人市场 2025 年已达数十亿美元规模，并将以约 13\%--17\% 的年复合增长率增长至百亿美元以上（表~\ref{tab:market-forecast}），其主要驱动力正是石化、电力等高危场景的自动化巡检需求。

\begin{table}[htbp]
\centering
\caption{全球巡检机器人市场规模预测}
\label{tab:market-forecast}
\begin{tabularx}{\textwidth}{Xccc}
\toprule
\textbf{研究机构} & \textbf{2025 年规模} & \textbf{远期规模} & \textbf{CAGR} \\
\midrule
Market Research Future & 47.53 亿美元 & 225.4 亿美元（2035） & 16.84\% \\
Future Market Insights & 32 亿美元 & 117 亿美元（2035） & 13.9\% \\
Maximize Market Research & 67.2 亿美元 & 161.6 亿美元（2032） & 13.2\% \\
\bottomrule
\end{tabularx}
\end{table}

\section{政策与技术背景}
\label{sec:policy-background}

本项目的技术路线与国家产业政策和近年技术成熟度直接相关。

政策层面，与本项目最相关的有三点：（1）2021 年工信部等 15 部门印发的《“十四五”机器人产业发展规划》将安防巡检、危险环境作业类机器人列为重点方向；（2）2025 年国务院《政府工作报告》首次写入“具身智能”与“智能机器人”，北京、深圳、杭州等地随后出台三年行动方案；（3）信创替代持续推进，相关政策提出 2027 年央企国企关键软硬件全面国产替代的目标，覆盖操作系统、中间件等领域。本项目的操作系统（开源鸿蒙）、AI 算力（华为昇腾）与开发语言（ArkTS）均为国产自主，正对应上述方向。

技术层面，三项条件近年趋于成熟，使本系统可行：其一，激光 SLAM、粒子滤波定位、A* 规划等导航算法已可在嵌入式平台实时运行，本项目的难点与价值在于不借助 ROS 生态、在 OpenHarmony 上自主实现完整算法栈；其二，以昇腾 310B 为代表的边缘 NPU 使 YOLO 检测、关键点回归等视觉模型可在机器人本体实时推理，无需云端；其三，DeepSeek-R1-Distill 等蒸馏小模型的出现，使大模型可在边缘 NPU 上离线运行，机器人本体可直接生成自然语言分析报告。这三项条件使“自主移动 + 视觉读数 + 智能分析”的全部能力可以完全由国产端侧软硬件承载。

\section{目标问题定义}
\label{sec:problem-definition}

综合上述背景，本项目要解决四个相互关联的问题：

\begin{enumerate}
\item 自主导航：在未知、无 GPS 的室内环境中，依靠激光雷达与轮控里程估计自主完成建图与定位，并在已建地图上实现全局定位与避障导航，且不依赖 ROS 及任何境外机器人中间件。
\item 全覆盖巡检：在含非凸形状与内部障碍的栅格地图上，规划一条完备、低重复的全覆盖路径；并能把一片大区域均衡划分给多台机器人协同覆盖。
\item 仪表读数：在边缘 NPU 上实时检测工业仪表、定位表盘关键点，通过几何换算得到指针读数，支持一个画面内多块仪表同时识别与越限告警。
\item 多机协同：不依赖中心服务器，使多台设备在同一可信网络中共享一致的任务状态，并将全局态势实时呈现给管理者。
\end{enumerate}

后续各章的设计与实现均围绕这四个问题展开。

\section{意义与价值}
\label{sec:significance}

在技术层面，本项目验证了国产操作系统承载复杂机器人系统的可行性：在 OpenHarmony 5.0 上不借助 ROS，自主实现雷达驱动、SLAM、规划与控制的完整闭环；同时验证了在鸿蒙分布式软总线上用共享任务状态组织多机协同、以及在同一块昇腾 NPU 上让实时视觉推理与端侧大模型分时共存的工程方案。

在应用层面，全栈国产自主的路线使系统可直接服务于电力、能源等对自主可控有刚性要求的场景；自主全覆盖巡检与实时读数可替代高危、低效的人工抄表作业；模块化的架构也便于扩展红外测温、气体检测等更多巡检模态，具有清晰的产业化前景。
